\documentclass[wabun2025.tex]{subfiles}
\begin{document}
\section{序論}
量子コンピュータが構想された当初から、量子化学シミュレーションは量子コンピュータ
によって大きく前進することが期待されている応用分野である
\cite{Cao2019.doi:10.1021/acs.chemrev.8b00803}。現時点で主流の量子コン
ピュータのハードウェアはノイズによる計算誤りがつきものであり、誤り訂正技術の実用
化が進められている。量子コンピュータ向けの誤り訂正技術として有望なsurface code法
では、古典コンピュータにおける誤り訂正技術に比べて高い冗長度を必要とする
\cite{Fowler2012.PhysRevA.86.032324}。そのため、誤り訂正を備えた実用的な量子コン
ピュータ（FTQC）の実現にはハードウェアデバイスにおけるqubit数の増大をはじめとす
る実装上の課題が存在し、実現にはまだ何年も掛かるとされている。誤り訂正が実装され
ておらず、量子ビット数も限られた現状のNISQ世代の量子コンピュータにおいては、量子
化学計算アルゴリズムとして第二量子化に基づくVQEベースのアルゴリズムが研究の主流
である\cite{peruzzo2014variational}。VQEでは量子コンピュータと古典コンピュータを
併用して必要な量子ビット数と量子回路の深さを抑えることができ、ノイズの存在を許容
できるので、NISQでも実行が可能である。しかし、FTQCが実現されることでノイズの影響
を抑えることができ、さらに量子ビット数の制約が緩和されれば、より広い範囲のアルゴ
リズムが研究・実用化されることが期待される。

第一量子化に基づく量子化学シミュレーションの代表例であるグリッド法は、基底関数の
選択に依存しないため電子構造の変化を扱うことができ、また電子間の相互作用を直接扱
えるため、量子化学シミュレーションとしては究極的なターゲットの一つであるが、必要
とする計算資源はVQEに比べて多い\cite{kassal2008.pnas.0808245105}。グリッド法では
等間隔の格子点のインデクスを複数のqubitからなるレジスタの計算基底（のテンソル
積）に対応付け、その基底の振幅に各格子点の波動関数の振幅を格納し、波動関数の時間
発展を計算する。時間発展の計算アルゴリズムとしてはSuzuki-Trotter分解に基づくアル
ゴリズムが広く知られているが、より洗練されたアプローチによるアルゴリズムも提案さ
れている\cite{Berry2015.PhysRevLett.114.090502,
Babbush_2018,YuanSu.PRXQuantum.2.040332, eklund2026.arXiv2603.19007}。しかし、グ
リッド法のアルゴリズムはいずれも必要とするqubit数と回路の深さが大きく、かつ計算
エラーを許容できないため、アルゴリズムの動作検証の報告例は限られている。著者の知
る限り、量子コンピュータ実機を用いた計算例はまだ報告されておらず、既存の報告例は
いずれも古典計算機上の量子コンピュータエミュレータを用いた計算例である。例え
ば、Benentiらは、簡略化されたポテンシャル関数を用いた計算を行っている
\cite{Benenti2008.10.1119/1.2894532}。Chanらは量子コンピュータエミュレータと古典
プログラムを併用して、クーロンポテンシャルの計算は古典プログラム側で処理している
\cite{Chan2023.sciadv.abo7484}。すなわち、クーロンポテンシャルの計算自体を完全に
量子回路で実現した例は、エミュレータ上といえども報告されていない。そこで本研究で
は、クーロンポテンシャルの計算を含めて全て量子回路で実装したグリッド法のシミュ
レータの構成と動作例を示す。シミュレータ回路を生成するソフトウェアには、著者が開
発したcrsQを用いる\cite{takahashi2024.doi:10.1021/acs.jctc.4c00708}。グリッド法
による波動関数の時間発展計算を全て量子回路で実装し、量子コンピュータエミュレータ
上で動作させ、計算結果の評価を試みる。そのために、以下の3つの課題に取り組む。

一つ目の課題は、限られたqubit数のエミュレータで動作可能な回路構成を示すことであ
る。本研究ではエミュレータに収まるように計算モデルを単純化し、原子核座標は固定し
た単一水素原子の1次元と2次元のモデルを対象とする。クーロンポテンシャルを計算する
基本的な方法は二進の算術演算に基づく方法である\cite{Zalka-rspa.1998.0162}。この
方法に対して、計算の一部を古典計算機上で済ませておきQROMを用いてルックアップする
ことによりゲート数を削減する方法が提案されているが、補助qubitを多く必要とする
\cite{Sanders2020.PRXQuantum.1.020312}。本研究では、 qubit数を節約するため、制御
レジスタの値に応じてターゲットqubitを回転させることができるUCRZ回路
\cite{Mottonen2004.PhysRevLett.93.130502}を採用する。時間発展計算には基本的な
Suzuki-Trotter分解を用いる\cite{Zalka-rspa.1998.0162}。これによってエミュレータ
や初期のFTQCで動作可能な回路構成例を提示する。グリッド法においては、波動関数の初
期化回路の規模も課題となり、任意の状態準備にはグリッドサイズの指数に対するゲート
数を必要とする。この規模を削減する方法として、原子軌道の線形結合で表現された初期
状態をMatrix ProductState(MPS)形式(あるいはTensor Train形式
\cite{Oseledets2011.doi:10.1137/090752286})に変換し、低ランク近似に基づいて初期
化回路を構成する方法が提案されている\cite{Huggins2025.PRXQuantum.6.020319}。この
手法は、波動関数の滑らかさや粒子間の相関の強さに応じたbond rankを用いることで初
期化回路の規模を大幅に削減可能とするものである。しかし、本研究ではエミュレータで
の実行が可能な小規模なグリッドサイズを対象とするため、MPSによる圧縮は行わず、実
装が容易な単純な初期化回路を用いる。

二つ目の課題は、クーロンポテンシャルの特異点の扱いを定めることである。グリッド法
では粒子座標の組合せに対してポテンシャルエネルギーを計算し、その値に比例した角度
で波動関数の位相を更新する必要がある。原子核と電子の座標が重なる点はクーロンポテ
ンシャルの特異点であり、数値計算上は発散する。波動関数を展開する基底として平面波
基底（運動量空間表現）を用いる方法では、ポテンシャルを運動量空間での演算として扱
うことで、実空間の特異点の直接的な計算を回避できることが知られている
\cite{Babbush_2019, eklund2026.arXiv2603.19007}。ただしこれらの方式は補助qubitを
多く必要とし、小規模な量子コンピュータでの実行は望めない。一方で位置基底を用いる
アプローチは、補助qubit数の少ない回路で実現できるが、特異点への対処が不可欠とな
る。これに対し、粒子間距離にオフセットを加えるソフトコアポテンシャルを用いる方式
\cite{Kivlichan_2017}や、原子核の格子と電子の格子の定義に相対的なズレを与えて距
離がゼロになるのを避ける方式\cite{Chan2023.sciadv.abo7484}などが提案されている。
本研究では、特異点を含む格子区画において、ポテンシャル関数の当該領域における積分
平均値（あるいは体積平均）に基づいて実効的なポテンシャル値を定める方法を提案し、
数値評価によりその妥当性を検証する。

三つ目の課題は、与えられたqubit数の中で誤差を最小化するシミュレーションパラメー
タの決定方法を確立することである。グリッド法ではシミュレーションパラメータとし
て、まず空間の分解能（格子間隔）とシミュレーション領域の大きさを定める必要があ
る。さらに、Suzuki-Trotter分解を用いる場合には時間発展の刻み幅も定める必要があ
る。Suzuki-Trotter分解では本来非可換であるハミルトニアンの位置エネルギー項と運動
エネルギー項を含むハミルトニアンの指数関数を、それぞれの指数関数の積へと分解す
る。この分解に伴い、演算子の非可換性に起因する誤差が生じる。初期の誤差分析は、演
算子ノルムに基づき、あらゆる入力状態に対する最悪値を評価するものであった。そのた
め、現実の数値計算で経験的に知られている誤差よりも数桁過大に見積もられることが指
摘されていたが、演算子の交換子スケーリングに着目することで、よりタイトな誤差評価
が可能となった\cite{Childs2021.PhysRevX.11.011020}。その後、最悪値評価の代わり
に、入力状態がハミルトニアンの定義域に対して持つ特性、すなわち状態依存性
(state-dependency)を考慮することで、評価の余裕をさらに縮小できることが示された
\cite{Burgarth2024.PhysRevResearch.6.043155}。同研究ではポテンシャルの特異点にお
いて波動関数の値が非ゼロの場合($s$軌道など)、誤差の収束がトロッター分割数$n$に関
して$n^{-1/4}$という劣線形(sublinear)なスケーリングに留まることや、高次トロッ
ター公式がこのような状態に対しては精度向上に寄与しないことも示されている。この結
果を多粒子に拡張した研究ではクーロンポテンシャルのような非有界ハミルトニアン
(unbounded Hamiltonian)において、分割数$n$への依存性が$n^{-1/4}$であり、粒子数に
対する誤差の増大が多項式的であることが証明されている\cite{fang2026trotter}。本研
究では、エミュレータと古典プログラムに基づく数値評価により、単一電子のモデルにお
いて空間離散化誤差と領域打ち切り誤差のトレードオフから空間パラメータを最適化し、
サンプリング定理の要請と最近のトロッター誤差評価を考慮して時間パラメータを決定す
る統合的なアプローチを提示する。

\end{document}
